\documentclass[../../main.tex]{subfiles}
\begin{document}

% 年度の大きな表示
\begin{center}
  {\fontsize{48pt}{58pt}\selectfont \textbf{2005}\normalsize\textbf{年度}}
\end{center}

\vspace{10mm}

% 本文


\vspace{15mm}

% 出題テーマと難易度の表
\noindent\textbf{出題テーマと難易度}

\vspace{5mm}

\begin{center}
  \shadowbox{%
    \begin{tabular}{|c|p{8cm}|c|}
      \hline
      \textbf{問題番号} & \textbf{テーマ} & \textbf{難易度} \\
      \hline
      第1問 & 定積分の漸化式と不等式評価（$(\log x)^n$の積分） & \\
      \hline
      第2問 & 確率と期待値（さいころ2回・多項式の係数） & \\
      \hline
      第3問 & 空間図形（円板が通過する部分の体積） & \\
      \hline
      第4問 & 円の媒介変数表示と対称式の存在範囲・最大最小 & \\
      \hline
    \end{tabular}%
  }
\end{center}

\vspace{15mm}

% 解答
\noindent\textbf{解答}

\vspace{5mm}

\begin{center}
  \shadowbox{%
    \renewcommand{\arraystretch}{1.6}
    \begin{tabular}{|c|c|p{8cm}|}
      \hline
      \textbf{問題番号} & \textbf{小問} & \textbf{答え} \\
      \hline
      \multirow{3}{*}{第1問} & (1) & 証明問題 \\
      \cline{2-3}
                             & (2) & 証明問題 \\
      \cline{2-3}
                             & (3) & 証明問題 \\
      \hline
      \multirow{2}{*}{第2問} & (1) & $E(s)=-7$，$E(t)=\dfrac{49}{4}$ \\
      \cline{2-3}
                             & (2) & $E(a)=-14$，$E(b)=\dfrac{238}{3}$，$E(c)=-\dfrac{637}{3}$，$E(d)=\dfrac{8281}{36}$ \\
      \hline
      第3問 & - & $V=2\pi+\dfrac{16}{3}$ \\
      \hline
      \multirow{2}{*}{第4問} & (1) & 図示（$\dfrac{s^2-1}{2}\le t\le\dfrac{s^2}{4}$，$-\sqrt{2}\le s\le\sqrt{2}$） \\
      \cline{2-3}
                             & (2) & 最大値$\sqrt{2}m+\dfrac{1}{2}$；最小値は$0\le m\le\sqrt{2}$で$-\dfrac{1}{2}m^2-\dfrac{1}{2}$，$m\ge\sqrt{2}$で$-\sqrt{2}m+\dfrac{1}{2}$ \\
      \hline
    \end{tabular}%
    \renewcommand{\arraystretch}{1.0}
  }
\end{center}

\vspace{15mm}

% 解答時間
\noindent\textbf{試験形態}

\vspace{5mm}

\begin{center}
  \shadowbox{%
    \begin{tabular}{p{8cm}}
      （要確認）
    \end{tabular}%
  }
\end{center}

\end{document}
